\section{Methodology}


%================================================
% \begin{frame}{Methodology Overview}

% 	\centering

% 	\begin{tikzpicture}[
% 			node distance=0.4cm,
% 			every node/.style={
% 					draw,
% 					rounded corners,
% 					align=center,
% 					minimum width=3.8cm,
% 					minimum height=0.8cm
% 				},
% 			>=Stealth
% 		]

% 		%-----------------------------------
% 		% Experimental branch

% 		\node (video)
% 		at (-4,0)
% 		{Experimental Videos};

% 		\node (track)
% 		at (-4,-1.5)
% 		{Detection \& Tracking};

% 		\node (traj)
% 		at (-4,-3)
% 		{Droplet Trajectories};

% 		%-----------------------------------
% 		% Model branch

% 		\node (model)
% 		at (4,0)
% 		[fill=red!10,draw=red!70!black]
% 		{Memory-Driven Model};

% 		\node (pde)
% 		at (4,-1.5)
% 		{PDE--SDE Formulation};

% 		\node (sim)
% 		at (4,-3)
% 		{Numerical Simulation};

% 		%-----------------------------------
% 		% Comparison

% 		\node (compare)
% 		at (0,-4.8)
% 		[fill=red!15,draw=red!70!black]
% 		{Experiment--Simulation Comparison};

% 		%-----------------------------------
% 		% Arrows

% 		\draw[->,thick] (video) -- (track);
% 		\draw[->,thick] (track) -- (traj);

% 		\draw[->,thick] (model) -- (pde);
% 		\draw[->,thick] (pde) -- (sim);

% 		\draw[->,thick] (traj) -- (compare);
% 		\draw[->,thick] (sim) -- (compare);

% 	\end{tikzpicture}
% \end{frame}
%================================================
\begin{frame}{Experimental System}

	\begin{columns}

		\column{0.55\textwidth}

		\begin{itemize}
			\item Self-propelled Belousov--Zhabotinsky (BZ) droplets in a monoolein--squalane oil medium.
			\item Droplet diameter: $\approx 1.75\,\mathrm{mm}$.
			\item Videos recorded using a CMOS camera at 5 fps.
			\item Spatial resolution: $0.06\,\mathrm{mm}$ per pixel.

		\end{itemize}

		\column{0.45\textwidth}

		\centering

		% Replace with experimental setup image
		\includegraphics[height=0.58\textheight]{camera_setup_image.png}

		\vspace{0.2cm}
		{\scriptsize Experimental setup for recording active droplet trajectories. Adapted from Meshram and Shajahan (2026).\par}
	\end{columns}
	\rightcite{D5SM01187F}

\end{frame}

%================================================
%================================================
\begin{frame}{Trajectory Extraction}

	\vspace{-0.2cm}

	\begin{columns}[T]

		%=========================================
		\column{0.48\textwidth}

		\begin{block}{YOLOv8 Detection and Tracking}

			\small

			\begin{itemize}
				\item Detect droplets in each frame.
				\item Assign persistent tracking IDs.
				\item Link detections across frames.
				\item Extract centroid coordinates.
				\item Reconstruct droplet trajectories.
			\end{itemize}

			\vspace{0.15cm}

			{\footnotesize
				Transfer learning using YOLOv8
				(369 training images, 96 validation images,
				50 epochs, resolution $640\times640$).
			}

		\end{block}

		\vspace{0.15cm}

		%=========================================
		\column{0.52\textwidth}

		\centering

		\includegraphics[height=0.55\textheight]{yolo_tracking_example.png}

		\vspace{0.1cm}

		{\tiny
			\begin{minipage}{0.9\linewidth}
				\centering
				YOLOv8 detection and tracking output.\\
				Bounding boxes and persistent tracking IDs\\
				enable trajectory reconstruction across video frames.
			\end{minipage}
		}
	\end{columns}

	\begin{alertblock}{Output}

		Tracked droplet trajectories used for
		transport and memory analysis.

	\end{alertblock}
\end{frame}
%================================================
\begin{frame}{From Trajectories to Memory Measures}

	\centering

	\begin{tikzpicture}[
			node distance=0.6cm,
			every node/.style={
					draw,
					rounded corners,
					align=center,
					minimum width=3cm,
					minimum height=1cm
				},
			> =Stealth
		]

		\node (traj) {Trajectory Data};

		\node (msd) [right=of traj]
		{MSD};

		\node (vacf) [right=of msd]
		{VACF};

		\node (memory) [right=of vacf,
			fill=red!15,
			draw=red!70!black]
		{Memory};

		\draw[->,thick] (traj) -- (msd);
		\draw[->,thick] (msd) -- (vacf);
		\draw[->,thick] (vacf) -- (memory);

	\end{tikzpicture}

	\vspace{0.6cm}

	\begin{columns}

		\column{0.4\textwidth}

		\begin{block}{Transport Measures}
			\begin{itemize}
				\item Mean Squared Displacement (MSD)
				\item Local MSD Exponent
			\end{itemize}
		\end{block}

		\column{0.4\textwidth}

		\begin{alertblock}{Memory Measures}
			\begin{itemize}
				\item Velocity Autocorrelation Function (VACF)
				\item Orientation Correlation Function (OCF)
			\end{itemize}
		\end{alertblock}

	\end{columns}

\end{frame}
\begin{frame}{Transport Characterization}

	\begin{columns}[T]

		%----------------------------------
		\column{0.48\textwidth}

		\begin{block}{Mean Squared Displacement (MSD)}

			\vspace{-0.05cm}

			\[
				\mathrm{MSD}(\tau)
				=
				\left\langle
				\|\mathbf{X}(t+\tau)-\mathbf{X}(t)\|^2
				\right\rangle
			\]

		\end{block}

		\begin{itemize}
			\item Measures spatial exploration.
			\item Identifies ballistic, diffusive, and anomalous transport.
			\item Used to estimate transport behavior.
		\end{itemize}

		%----------------------------------
		\column{0.48\textwidth}

		\begin{block}{Local MSD Exponent}

			\vspace{-0.05cm}

			\[
				\alpha(\tau)
				=
				\frac{d\log[\mathrm{MSD}(\tau)]}
				{d\log(\tau)}
			\]

		\end{block}

		\begin{itemize}
			\item Characterizes lag-dependent transport.
			\item $\alpha \approx 2$ : Ballistic motion
			\item $\alpha \approx 1$ : Normal diffusion
			\item $1<\alpha<2$ : Superdiffusion
		\end{itemize}

	\end{columns}

	\vspace{0.2cm}

	\begin{alertblock}{Physical Interpretation}
		MSD quantifies how far droplets travel, while the local exponent reveals how the transport regime evolves across timescales.
	\end{alertblock}

\end{frame}
\begin{frame}{Memory Characterization}

	\begin{columns}[T]

		%----------------------------------
		\column{0.4\textwidth}
		\vspace{-0.2cm}

		\begin{alertblock}{Velocity Autocorrelation Function (VACF)}

			\[
				C_v(\tau)
				=
				\frac{
					\left\langle
					\mathbf{v}(t)\cdot\mathbf{v}(t+\tau)
					\right\rangle}
				{
					\left\langle
					|\mathbf{v}(t)|^2
					\right\rangle}
			\]
		\end{alertblock}

		\vspace{0.05cm}

		\[
			C_v(\tau)\sim e^{-\tau/\tau_p}
		\]

		\begin{itemize}
			\item Measures velocity memory.
			\item Provides persistence time $\tau_p$.
		\end{itemize}

		%----------------------------------
		\column{0.4\textwidth}
		\vspace{-0.2cm}

		\begin{alertblock}{Orientation Correlation Function (OCF)}
			\vspace{-0.2cm}

			\[
				C_\theta(\tau)
				=
				\left\langle
				\cos[\theta(t+\tau)-\theta(t)]
				\right\rangle
			\]
			% \vspace{0.00005cm}
			% \vspace{-0.1cm}

		\end{alertblock}

		\vspace{0.05cm}

		\[
			C_\theta(\tau)\sim e^{-\tau/\tau_r}
		\]

		\begin{itemize}
			\item Measures directional memory.
			\item Provides rotational persistence time $\tau_r$.
		\end{itemize}

	\end{columns}

	\vspace{0.3cm}

	\begin{alertblock}{Physical Meaning}
		These observables quantify how long information about past velocity and direction persists.
	\end{alertblock}

\end{frame}
